\documentclass[aspectratio=169,10pt]{beamer}
% Shared Beamer setup for Fall 2026 course slides.
% Clean Cal Poly-inspired theme: no custom class files or uncommon packages.
\mode<presentation>{
  \usetheme{default}
  \usefonttheme{professionalfonts}
}

\definecolor{CalPolyGreen}{HTML}{154734}
\definecolor{CalPolyGold}{HTML}{BD8B13}
\definecolor{CalPolySage}{HTML}{7FA28F}
\definecolor{CalPolyMint}{HTML}{DDF1E7}
\definecolor{CalPolySand}{HTML}{A29F86}
\definecolor{CalPolyBeige}{HTML}{EFEEDE}
\definecolor{DeckInk}{HTML}{1F2933}
\definecolor{DeckMuted}{HTML}{5F6B7A}
\definecolor{DeckSoft}{HTML}{F7F7F5}

\setbeamersize{text margin left=0.55in,text margin right=0.55in}
\setbeamercolor{normal text}{fg=DeckInk,bg=white}
\setbeamercolor{structure}{fg=CalPolyGreen}
\setbeamercolor{title}{fg=white}
\setbeamercolor{frametitle}{fg=CalPolyGreen,bg=white}
\setbeamercolor{block title}{bg=CalPolySage,fg=white}
\setbeamercolor{block body}{bg=CalPolyMint,fg=black}
\setbeamercolor{alerted text}{fg=CalPolyGold}

\setbeamerfont{title}{size=\Huge,series=\bfseries}
\setbeamerfont{frametitle}{size=\Large,series=\bfseries}
\setbeamerfont{block title}{size=\normalsize,series=\bfseries}

\setbeamertemplate{navigation symbols}{}
\setbeamertemplate{itemize item}{\color{CalPolyGold}$\blacktriangleright$}
\setbeamertemplate{itemize subitem}{\color{CalPolyGreen}$\bullet$}
\setbeamertemplate{footline}{
  \hfill{\color{DeckMuted}\scriptsize\insertframenumber/\inserttotalframenumber}
  \hspace{0.18in}\vspace{0.12in}
}
\setbeamertemplate{frametitle}{
  \vspace{0.18in}
  \hspace{0.02in}\insertframetitle\par
  \vspace{0.08in}
  \noindent\makebox[\linewidth][l]{%
    {\color{CalPolyGold}\rule{0.55in}{2pt}}%
    {\color{CalPolyGreen}\rule{\dimexpr\linewidth-0.55in\relax}{2pt}}%
  }
}

\newcommand{\CourseNumber}{Course}
\newcommand{\CourseTitle}{Course Title}
\newcommand{\LectureNumber}{0}
\newcommand{\LectureTitle}{Lecture Title}
\newcommand{\LectureDate}{Date}
\newcommand{\CourseUnit}{Unit}

\newcommand{\coursenumber}[1]{\renewcommand{\CourseNumber}{#1}}
\newcommand{\coursetitle}[1]{\renewcommand{\CourseTitle}{#1}}
\newcommand{\lecturenumber}[1]{\renewcommand{\LectureNumber}{#1}}
\newcommand{\lecturetitle}[1]{\renewcommand{\LectureTitle}{#1}}
\newcommand{\lecturedate}[1]{\renewcommand{\LectureDate}{#1}}
\newcommand{\courseunit}[1]{\renewcommand{\CourseUnit}{#1}}

\author{Kirk Duran}
\institute{California Polytechnic State University, San Luis Obispo}
\date{\LectureDate}

\newcommand{\makedecktitle}{
  \begingroup
  \setbeamercolor{background canvas}{bg=CalPolyGreen}
  \begin{frame}[plain]
    \vfill
    \begin{center}
      {\usebeamerfont{title}\color{white}\LectureTitle\par}
      \vspace{0.18in}
      {\Large\color{white}Lecture \LectureNumber\par}
    \end{center}
    \vfill
    \noindent{\color{white}\rule{\linewidth}{1.2pt}}\par
    \vspace{0.08in}
    \noindent\colorbox{CalPolyGold}{\makebox[0.24\linewidth][c]{\color{white}\Large\CourseNumber}}
    \hspace{0.12in}{\color{white}\Large Kirk Duran}
  \end{frame}
  \endgroup
}

\newenvironment{keyidea}{\begin{block}{Key idea}}{\end{block}}
\newenvironment{activity}{\begin{block}{In class}}{\end{block}}
\newenvironment{cpdefinition}{\begin{block}{Definition}}{\end{block}}
\newenvironment{exampleblockcp}{
  \begingroup
  \setbeamercolor{block title}{bg=CalPolySand,fg=white}
  \setbeamercolor{block body}{bg=CalPolyBeige,fg=black}
  \begin{block}{Example}
}{
  \end{block}
  \endgroup
}

\newcommand{\imageplaceholder}[2][0.3in]{%
  \noindent\makebox[\linewidth][c]{%
    \fcolorbox{CalPolySage}{CalPolyBeige}{%
      \parbox[c][#1][c]{0.86\linewidth}{%
        \centering
        {\color{DeckMuted}\scriptsize Image placeholder: }%
        {\color{DeckInk}\scriptsize\ttfamily\detokenize{#2}}%
      }%
    }%
  }%
}

\newcommand{\sectionframe}[1]{
  \begingroup
  \setbeamercolor{background canvas}{bg=CalPolyGreen}
  \begin{frame}[plain]
    \vfill
    {\color{white}\LARGE\bfseries #1\par}
    \vspace{0.18in}
    {\color{CalPolyGold}\rule{0.22\paperwidth}{3pt}}
    {\color{white}\rule{0.62\paperwidth}{3pt}}
    \vfill
  \end{frame}
  \endgroup
}

\newcommand{\placeholdernote}{
  \begin{block}{Placeholder}
    Replace these bullets with the final lecture material, examples, and in-class activities.
  \end{block}
}

\usepackage[T1]{fontenc}
\usepackage{lmodern}
\usepackage{amsmath,amssymb}
\usepackage{booktabs}
\pdfstringdefDisableCommands{\def\translate#1{#1}}
\graphicspath{{../assets/lecture-27/}}

% Lecture 27 follows the Lecture 26 architecture:
% Cal Poly title, CalPolyGreen concept leads, pale calculation boxes,
% keyidea callouts, and charts/figures alongside brief technical prose.
% Figures are regenerated with ../assets/lecture-27/generate-charts.py.
% The conceptual questioning illustration uses built-in image generation.
% Tree figures show the actual fitted models' top two levels for legibility;
% full model diagrams are also saved in the same asset directory.
% Primary references: https://scikit-learn.org/stable/modules/tree.html
% and https://scikit-learn.org/stable/modules/generated/sklearn.tree.plot_tree.html
% Worked printer data and synthetic evaluation differ from the robot lab.
% Audience: senior-level AI students without advanced mathematics prerequisites.
% Lead with English; equations are supporting notation, not the entry point.
% Estimated pacing: motivation/tree language 10 min; entropy & gain 12;
% recursive construction 8; numeric thresholds 5; Python demos 6;
% overfitting/pruning/validation 7; synthesis 2. Total 50 minutes.
% All content is in the lecture; no ensemble methods or appendix slides.

\newcommand{\lecturevisual}[3][2.15in]{%
  \par\vspace{2pt}%
  {\centering\includegraphics[width=.98\linewidth,height=#1,keepaspectratio]{#2}\par}%
  \vspace{2pt}%
}
\newcommand{\conceptlead}[1]{%
  {\normalsize\bfseries\color{CalPolyGreen}#1\par}%
}
\newenvironment{calculation}[1][How we calculate it]{%
  \par\begingroup
  \setbeamercolor{calculation body}{bg=DeckSoft,fg=DeckInk}%
  \begin{beamercolorbox}[sep=2pt,wd=\linewidth]{calculation body}%
  \footnotesize{\bfseries\color{DeckMuted}#1\par}%
  \setlength{\abovedisplayskip}{3pt}%
  \setlength{\belowdisplayskip}{3pt}%
}{\end{beamercolorbox}\endgroup}
\setbeamersize{text margin left=.42in,text margin right=.42in}
\coursenumber{CS 4880}
\coursetitle{Artificial Intelligence}
\lecturenumber{27}
\lecturetitle{Decision Trees: Recursive Partitioning, Information Gain, and Pruning}
\lecturedate{October 28, 2026}
\courseunit{Learning Agents and Machine Learning}

\begin{document}
\makedecktitle

\section{Why Decision Trees?}

\begin{frame}{From Smooth Fits to Explicit Decisions}
\small
\conceptlead{Sometimes the world behaves more like a collection of rules than one smooth trend.}
\begin{columns}[T,onlytextwidth]
 \begin{column}{.49\textwidth}
  \begin{itemize}
   \item Lecture 26 fit lines, curves, and logistic scores to labeled outcomes.
   \item Printer failure risk may change abruptly above a job-load threshold; temperature may change the rule for lighter jobs.
   \item A tree learns \textbf{which question to ask next} and makes separate predictions in different regions.
  \end{itemize}
  \begin{keyidea}A tree can model sharp thresholds and interactions without adding polynomial features by hand.\end{keyidea}
 \end{column}
 \begin{column}{.48\textwidth}
   \lecturevisual[2.18in]{smooth_vs_rule.png}{Smooth prediction contrasted with an explicit threshold.}
 \end{column}
\end{columns}
\end{frame}

\begin{frame}{Akinator: The Value of the Next Question}
\small
\conceptlead{An informative question separates the remaining possibilities into useful groups.}
\begin{columns}[T,onlytextwidth]
 \begin{column}{.48\textwidth}
  \begin{itemize}
   \item Akinator provides an intuition: some questions eliminate more possibilities than others.
   \item ``Is it fictional?'' can be more useful than ``Does it wear a hat?'' depending on the candidates still remaining.
   \item A decision tree learns feature questions from labeled examples. The illustration is an analogy, not a description of Akinator's implementation.
  \end{itemize}
 \end{column}
 \begin{column}{.48\textwidth}
  \lecturevisual[1.9in]{akinator_questions.png}{An original conceptual illustration of informative questioning.}
 \end{column}
\end{columns}
\end{frame}

\begin{frame}{Anatomy of a Decision Tree}
\small
\conceptlead{Prediction means following answers from the root until a leaf supplies a label.}
\begin{columns}[T,onlytextwidth]
 \begin{column}{.51\textwidth}
  \begin{itemize}
   \item \textbf{Root:} the first test applied to every example.
   \item \textbf{Internal node:} another feature-based question.
   \item \textbf{Branch:} one answer to that question.
   \item \textbf{Leaf:} final prediction; for classification, usually the local majority class.
  \end{itemize}
  One root-to-leaf path is an \textbf{if--then rule}. The split question and predicted label play different roles.
 \end{column}
 \begin{column}{.46\textwidth}
  \lecturevisual[2.2in]{tree_regions.png}{The tree and the input regions it defines.}
 \end{column}
\end{columns}
\end{frame}

\begin{frame}{Recursive Partitioning Creates Decision Boundaries}
\small
\conceptlead{Every test separates the observations; later tests refine only one resulting group.}
\begin{itemize}
 \item \textbf{Categorical feature:} ``Warm = Yes?'' produces groups based on feature values.
 \item \textbf{Numeric feature:} ``Job size $\leq 6$?'' puts values on opposite sides of a threshold.
 \item Multiple tests create rectangular, \emph{axis-aligned} regions when plotted against two numerical features.
 \item Interactions arise from \emph{the path}: the temperature question can matter only after job size was tested.
\end{itemize}
\begin{calculation}[One path as a compact rule]
\[
(\text{Job size}\leq 6)\ \land\ (\text{Temperature}\leq 70)
\quad\Longrightarrow\quad\widehat y=\text{Complete}.
\]
$\land$ means ``and.'' The tree represents this combination by following two branches, not by computing an explicit product of features.
\end{calculation}
\end{frame}

\section{Entropy and Information Gain}

\begin{frame}{What Makes One Split Better Than Another?}
\small
\conceptlead{A helpful question creates groups whose observed labels mostly agree.}
\begin{columns}[T,onlytextwidth]
 \begin{column}{.54\textwidth}
  \begin{itemize}
   \item A \textbf{pure group} contains one class only: no disagreement about its training labels.
   \item A \textbf{mixed group} contains both failed and complete observations.
   \item Many tiny branches can appear pure on training data without predicting well. Validation and pruning help judge whether the split is useful.
   \item Large and small children must not contribute equally: weight by how many examples each contains.
  \end{itemize}
 \end{column}
 \begin{column}{.42\textwidth}
  \begin{block}{Compare two possible questions}
  \textbf{Good:} 5 failed / 0 complete versus 1 failed / 4 complete.\par\medskip
  \textbf{Weak:} 4 failed / 1 complete versus 2 failed / 3 complete.
  \end{block}
 \end{column}
\end{columns}
\end{frame}

\begin{frame}{Entropy Measures Uncertainty in a Group}
\small
\conceptlead{Entropy is low when the answer is nearly certain and high when classes are balanced.}
\begin{columns}[T,onlytextwidth]
 \begin{column}{.51\textwidth}
   \begin{itemize}
    \item All failed or all complete: \textbf{zero entropy}.
    \item Half failed and half complete: \textbf{highest binary entropy}, one bit.
    \item A 75\% / 25\% split is uncertain, but less uncertain than 50\% / 50\%.
   \end{itemize}
   \begin{calculation}[Binary entropy: notation for this intuition]
   \[
   H(p)=-p\log_2(p)-(1-p)\log_2(1-p).
   \]
   $p$ is the fraction of failed examples; $1-p$ is the complete fraction. Treat $0\log_2 0$ as zero.
   \end{calculation}
 \end{column}
 \begin{column}{.46\textwidth}
  \lecturevisual[2.38in]{entropy_curve.png}{Binary entropy rises toward a balanced class distribution.}
 \end{column}
\end{columns}
\end{frame}

\begin{frame}{A Small Dataset We Can Trace by Hand}
\small
\conceptlead{Predict whether a printer job fails using two inputs recorded before starting.}
\begin{center}
\renewcommand{\arraystretch}{1.0}
\begin{tabular}{rlll}
\toprule
\textbf{Job} & \textbf{Load} & \textbf{Warm?} & \textbf{Observed outcome}\\\midrule
1 & High & Yes & Fail\\
2 & High & Yes & Fail\\
3 & High & Yes & Fail\\
4 & High & No & Fail\\
5 & High & No & Fail\\
6 & Low & Yes & Fail\\
7 & Low & Yes & Complete\\
8 & Low & No & Complete\\
9 & Low & No & Complete\\
10 & Low & No & Complete\\\bottomrule
\end{tabular}
\end{center}
The label is the outcome; Load and Warm are features. This synthetic example has six failures and four completions. Jobs 6 and 7 have identical recorded inputs.
\end{frame}

\begin{frame}{Calculate Uncertainty at the Root}
\small
\conceptlead{Start by counting labels. The ten printer jobs contain six failed and four complete.}
\begin{center}
\begin{tabular}{lccc}
\toprule
\textbf{Group} & \textbf{Failed (Y)} & \textbf{Complete (N)} & \textbf{Total}\\
\midrule
All observations & 6 & 4 & 10\\
\bottomrule
\end{tabular}
\end{center}
\begin{itemize}
 \item The failed fraction is $6/10$, so the group is mixed rather than pure.
 \item We need a score for this \emph{starting uncertainty} before evaluating questions.
\end{itemize}
\begin{calculation}[Substitute the two class fractions]
\[
H(S)=-\frac6{10}\log_2\!\left(\frac6{10}\right)
     -\frac4{10}\log_2\!\left(\frac4{10}\right)
     \approx \boxed{0.971\text{ bits}}.
\]
The logarithms convert class proportions into an uncertainty score; the main procedure is count $\rightarrow$ fraction $\rightarrow$ entropy.
\end{calculation}
\end{frame}

\begin{frame}{Information Gain: How Much Uncertainty Did the Question Remove?}
\small
\conceptlead{Subtract the uncertainty left after the split from the uncertainty we had before it.}
\begin{enumerate}
 \item Calculate the parent group's entropy.
 \item Split the records using a candidate question.
 \item Calculate entropy inside each child and weight by its share of the records.
 \item Prefer the question producing the \textbf{largest reduction}.
\end{enumerate}
\begin{calculation}[One formula for all candidate splits]
\[
\underbrace{IG(S,X)}_{\text{gain}}
 =\underbrace{H(S)}_{\text{before}}
 -\underbrace{\sum_{v}\frac{|S_v|}{|S|}H(S_v)}_{\text{weighted uncertainty after}}.
\]
$S_v$ is the group receiving answer $v$; $|S_v|/|S|$ is its fraction of observations. The sum simply adds the child contributions.
\end{calculation}
\end{frame}

\begin{frame}{Worked Split: Ask About Load or Warm?}
\small
\conceptlead{Load is the better root question: its children agree with their labels more strongly.}
\lecturevisual[1.4in]{split_comparison.png}{Exact class counts for the independent ten-job lecture example.}
\begin{calculation}[Compare both reductions from the same root entropy of 0.971 bits]
\[
\begin{aligned}
H_{\rm after}(\text{Load})&=\tfrac5{10}H(1)+\tfrac5{10}H(\tfrac15)=0.361, & IG&=0.610,\\
H_{\rm after}(\text{Warm})&=\tfrac5{10}H(\tfrac45)+\tfrac5{10}H(\tfrac25)=0.846, & IG&=0.125.
\end{aligned}
\]
Here $H(p)$ is binary entropy. The bigger gain leaves less uncertainty, so \textbf{Load becomes the root}.
\end{calculation}
\end{frame}

\section{Recursion and Stopping}

\begin{frame}{Build the Tree by Repeating the Same Decision}
\small
\conceptlead{Use the best local question, divide the records, then solve the smaller problems.}
\begin{columns}[T,onlytextwidth]
 \begin{column}{.55\textwidth}
 \begin{enumerate}
  \item \textbf{Stop} if labels are pure or a growth rule says to stop; make a leaf.
  \item \textbf{Search} candidate feature tests and compute each gain.
  \item \textbf{Choose} the largest useful gain and partition the data.
  \item \textbf{Recurse} separately on each child group.
 \end{enumerate}
 \begin{keyidea}``Greedy'' means best split \emph{now}. It does not guarantee the best complete tree; insisting on immediate gain can miss interactions that need another question.\end{keyidea}
 \end{column}
 \begin{column}{.42\textwidth}
 \begin{calculation}[Readable recursive pseudocode]
 \texttt{BuildTree(records):}\par
 \quad if stop: return majority leaf\par
 \quad find best candidate split\par
 \quad if no gain: return majority leaf\par
 \quad left = BuildTree(left records)\par
 \quad right = BuildTree(right records)\par
 \quad return node(test, left, right)
 \end{calculation}
 \end{column}
\end{columns}
\end{frame}

\begin{frame}{Continue the Example: A Small Complete Tree}
\small
\conceptlead{Only the mixed low-load branch needs another question.}
\lecturevisual[1.5in]{completed_tree.png}{Ask about warmth only in the low-load subgroup.}
\begin{itemize}
 \item \textbf{High load:} all five jobs failed; stop immediately.
 \item \textbf{Low load:} ask Warm. With No, all three complete; with Yes, one fails and one completes.
 \item Break the tied leaf using its parent's majority (complete). One training error remains because the recorded inputs cannot separate these jobs.
\end{itemize}
\end{frame}

\begin{frame}{Why Can the Tree Still Make a Mistake?}
\small
\conceptlead{A decision rule cannot distinguish observations that look identical to the model.}
\begin{block}{The inseparable pair in our data}
Jobs 6 and 7 both have \textbf{Load = Low} and \textbf{Warm = Yes}. One failed; the other completed.
\end{block}
\begin{itemize}
 \item Every deterministic tree using only these features sends both jobs down the same path.
 \item More splitting cannot recover a distinction that the inputs never recorded.
 \item Omitted sensors, random effects, noisy labels, or unrecorded conditions may explain the different outcomes.
\end{itemize}
\begin{calculation}[The limitation in compact form]
\[
\mathbf x_6=\mathbf x_7\quad\Longrightarrow\quad f(\mathbf x_6)=f(\mathbf x_7),
\quad\text{even though }y_6\neq y_7.
\]
\end{calculation}
\end{frame}

\section{Numerical Features}

\begin{frame}{Numerical Features: Where Should We Place the Cutoff?}
\small
\conceptlead{For a numerical feature, search possible ``less than or equal to'' cutoffs.}
\lecturevisual[1.3in]{thresholds.png}{Candidate thresholds for a separate six-job numerical example.}
\begin{enumerate}
 \item In a separate six-job example, sizes are $2,3,4,9,11,12$ ten-page batches.
 \item Candidate midpoints are $2.5,3.5,6.5,10,11.5$.
 \item For each midpoint, split the records and measure the impurity left behind.
 \item Choose the feature and threshold giving the greatest impurity reduction.
\end{enumerate}
\end{frame}

\begin{frame}{A Numeric Split Is the Same Impurity Calculation}
\small
\conceptlead{A threshold simply determines who goes left and right; the split score works as before.}
\begin{columns}[T,onlytextwidth]
 \begin{column}{.54\textwidth}
  \begin{itemize}
   \item Candidate: ``Job size $\leq 6.5$?''
   \item Place 2, 3, 4 on the left; 9, 11, 12 on the right.
   \item If every left observation is complete and every right observation is failed, both children are pure.
   \item Trying all midpoints is a search over possible decision rules, not advanced calculus.
  \end{itemize}
 \end{column}
 \begin{column}{.42\textwidth}
 \begin{calculation}[Weighted child impurity]
 \[
  I_{\rm split}(t)=\frac{n_L}{n}I(L)
   +\frac{n_R}{n}I(R).
 \]
 Smaller is better. Here each child is pure, so the score is zero.
 \end{calculation}
 \end{column}
\end{columns}
\begin{keyidea}The same mechanism works for classification (entropy or Gini) and regression (e.g., squared-error reduction).\end{keyidea}
\end{frame}

\begin{frame}{Stop Before a Tree Memorizes the Training Set}
\small
\conceptlead{Depth and leaf-size limits prevent the algorithm from inventing a rule for every tiny exception.}
\begin{columns}[T,onlytextwidth]
 \begin{column}{.49\textwidth}
  \begin{block}{\texttt{max\_depth}}
  Maximum number of splits along a path. Smaller values produce shallower, simpler trees.
  \end{block}
 \end{column}
 \begin{column}{.49\textwidth}
  \begin{block}{\texttt{min\_samples\_leaf}}
  Minimum number of training examples permitted in each leaf. Prevents tiny, fragile groups.
  \end{block}
 \end{column}
\end{columns}
\begin{itemize}
 \item These are \textbf{pre-pruning} controls: they constrain growth before the tree is finished.
 \item Too restrictive $\rightarrow$ meaningful patterns are missed (underfitting).
 \item Too flexible $\rightarrow$ sample-specific noise can become rules (overfitting).
\end{itemize}
\end{frame}

\section{Python: The Instructor's Two Examples}

\begin{frame}[fragile]{Python Example 1: California Housing Regression Tree}
\small
\conceptlead{The same recursive structure can predict a number rather than a class.}
{\scriptsize
\begin{verbatim}
import matplotlib.pyplot as plt
from sklearn.datasets import fetch_california_housing
from sklearn.tree import DecisionTreeRegressor, plot_tree

# Load data
data = fetch_california_housing()
X = data.data
y = data.target
feature_names = data.feature_names

# Fit for display; held-out evaluation comes later.
tree = DecisionTreeRegressor(max_depth=4, random_state=42)
tree.fit(X, y)

# Plot the tree
plt.figure(figsize=(20, 10))
plot_tree(tree, feature_names=feature_names,
          filled=True, rounded=True, fontsize=10)
plt.title("Decision Tree Regression on California Housing Dataset")
plt.show()
\end{verbatim}
}
\end{frame}

\begin{frame}{Reading the California Housing Tree}
\small
\conceptlead{A regression leaf predicts the mean target of its training records.}
\lecturevisual[1.4in]{california_tree.png}{Regenerated top two levels of the instructor's regression model; full tree saved with the assets.}
\begin{itemize}
 \item Read \texttt{MedInc} at the root: True goes left, False right. Top two levels shown; fitted model depth is four.
 \item \texttt{samples} counts records reaching the node; \texttt{value} is its mean target.
 \item \texttt{squared\_error} summarizes spread of targets within the node; reducing this spread replaces class entropy as the regression goal.
\end{itemize}
\end{frame}

\begin{frame}[fragile]{Python Example 2: Wine Classification Tree}
\small
\conceptlead{Now a leaf predicts a class using the observations that reach it.}
{\scriptsize
\begin{verbatim}
import matplotlib.pyplot as plt
from sklearn.datasets import load_wine
from sklearn.tree import DecisionTreeClassifier, plot_tree
# Load dataset
wine = load_wine()
X = wine.data
y = wine.target
feature_names = wine.feature_names
class_names = wine.target_names
# Fit for display; held-out evaluation comes later.
clf = DecisionTreeClassifier(max_depth=4, random_state=42)
clf.fit(X, y)
# Plot the tree
plt.figure(figsize=(20, 10))
plot_tree(clf, feature_names=feature_names,
          class_names=class_names, filled=True,
          rounded=True, fontsize=10)
plt.title("Decision Tree on Wine Dataset")
plt.show()
\end{verbatim}
}
\end{frame}

\begin{frame}{Reading the Wine Classification Tree}
\small
\conceptlead{Read each node's question, sample count, class distribution, and predicted class.}
\lecturevisual[1.5in]{wine_tree.png}{Regenerated top two levels of the instructor's wine classifier; full tree saved with the assets.}
\begin{itemize}
 \item \texttt{value} shows class counts; \texttt{class} is the majority. Two levels shown; fitted model depth is four.
 \item \texttt{gini} measures impurity; Gini is scikit-learn's default criterion.
 \item \texttt{criterion="entropy"} uses information gain instead. This model keeps the default Gini criterion.
\end{itemize}
\end{frame}

\section{Overfitting, Pruning, and Validation}

\begin{frame}{A Deep Tree Can Fit the Wrong Details}
\small
\conceptlead{Perfectly matching old observations does not mean predicting new jobs well.}
\begin{columns}[T,onlytextwidth]
 \begin{column}{.54\textwidth}
  \begin{itemize}
   \item Deep trees can carve out tiny regions to explain isolated training examples or misreadings.
   \item \textbf{Training error} often falls as complexity grows; \textbf{validation error} may stop improving or increase.
   \item Smaller trees have more bias (stronger simplifying assumptions); large trees can have more variance (sensitivity to the sample).
  \end{itemize}
 \end{column}
 \begin{column}{.43\textwidth}
  \lecturevisual[2.36in]{complexity_curve.png}{Measured training and validation errors on a separate synthetic classification dataset.}
 \end{column}
\end{columns}
\end{frame}

\begin{frame}{Two Ways to Prune: Before Growth or Afterward}
\small
\conceptlead{Pruning trades a little training fit for simpler rules that may generalize better.}
\begin{columns}[T,onlytextwidth]
 \begin{column}{.48\textwidth}
  \begin{block}{Pre-pruning}
  Stop growth using \texttt{max\_depth}, \texttt{min\_samples\_leaf}, or a minimum split improvement.
  \end{block}
  Easy to enforce, but may block useful later questions.
 \end{column}
 \begin{column}{.48\textwidth}
  \begin{block}{Post-pruning}
  Grow a larger tree, then remove branches when their improvement is too small to justify their complexity.
  \end{block}
  Cost-complexity pruning is controlled by \texttt{ccp\_alpha}.
 \end{column}
\end{columns}
\begin{calculation}[Cost-complexity pruning in one line]
\[
  R_\alpha(T)=R(T)+\alpha\,|\text{leaves}(T)|.
\]
$R(T)$ is the training impurity cost; $|\text{leaves}(T)|$ counts final leaves; $\alpha$ is a penalty for extra leaves. Larger $\alpha$ generally favors smaller trees.
\end{calculation}
\end{frame}

\begin{frame}{Choose the Tree Using Held-Out Evidence}
\small
\conceptlead{Compare models on held-out examples, not training accuracy alone.}
{\centering
\renewcommand{\arraystretch}{1.32}
\begin{tabular}{lccc}
\toprule
\textbf{Candidate} & \textbf{Leaves / depth} & \textbf{Training accuracy} & \textbf{Validation accuracy}\\
\midrule
Unrestricted & 50 / 10 & 100.0\% & 87.5\%\\
Pre-pruned & 9 / 4 & 88.6\% & 89.2\%\\
Post-pruned & 7 / 5 & 92.5\% & 92.5\%\\
\bottomrule
\end{tabular}
\par}
\begin{itemize}
 \item Separate synthetic experiment: 360 training, 120 validation, 120 test cases. Training fits rules; validation compares depth and pruning settings.
 \item Here the frozen post-pruned model scores 88.3\% on untouched test cases.
 \item If failures are rare, also examine missed failures and false alarms, not just accuracy.
\end{itemize}
\begin{keyidea}Removing a special-case leaf can make some training labels wrong but improve predictions on new jobs.\end{keyidea}
\end{frame}

\section{Interpretability and Takeaways}

\begin{frame}{Explain a Decision, Then Examine Its Failure}
\small
\conceptlead{The tree gives an inspectable chain of reasons; it cannot reveal what the sensors never measured.}
\begin{block}{Two readable rules from our printer tree}
\textbf{Rule 1:} IF load is high, THEN predict failed.\par\smallskip
\textbf{Rule 2:} IF load is low AND warm is no, THEN predict complete.
\end{block}
\begin{itemize}
 \item Trace a misclassification through its nodes: which test routed it to this leaf?
 \item Check the leaf's sample count: a rule based on one or two cases may be fragile.
 \item Missing features might include paper quality, component wear, or sensor confidence.
 \item Understandable rules help debugging, auditing, and explaining an agent's action to a human operator.
\end{itemize}
\end{frame}

\begin{frame}{Synthesis: From Questions to Generalization}
\small
\begin{enumerate}
 \item \textbf{Representation:} decision paths form conditional if--then rules and nonlinear, axis-aligned partitions.
 \item \textbf{Learning:} entropy describes class uncertainty; information gain chooses a useful local question; recursion builds the rest of the tree.
 \item \textbf{Limitations:} greedy choices are not globally optimal, and contradictory labels with identical features cannot be perfectly separated.
 \item \textbf{Generalization:} numeric thresholds, depth limits, leaf-size controls, pruning, and held-out validation determine whether those rules transfer.
\end{enumerate}
\begin{keyidea}A good decision tree asks informative questions \emph{and} knows when to stop asking.\end{keyidea}
\end{frame}

\end{document}
